\setcounter{section}{3}
\section{Inspect what aggregate scores hide}
\label{sec:episodes}

\subsection{Episode selection}
% Select ten RECORDED episodes: 3 typical successes, 3 typical failures,
% 2 surprising successes and 2 surprising failures. Justify typicality using
% the outcome distributions and explain the surprise criterion before interpretation.
% Record episode ID, source run/configuration, setting, agent and evaluation seed.
\noindent\textit{[Identify the ten episodes and justify their selection.]}

\subsection{Three important decisions per episode}
% Annotate THREE decisions in EACH of the ten selected episodes (30 in total).
% For each decision include: state and available actions; chosen action and actual
% Q-values or policy probabilities; observed transition and outcome; alternative
% action, what might have happened, and what remains uncertain.
% Use compact inline tables or trace figures with source run/episode/step IDs.
% Do not invent counterfactual transitions or probabilities absent from the traces.
\noindent\textit{[Insert the recorded decision annotations.]}

\subsection{What aggregate results miss}
% Use specific inspected episodes to distinguish poor decisions from unlucky
% transitions and explain at least one limitation of the aggregate scores.
\noindent\textit{[Explain the limitation using episode evidence.]}
